% Scope: analytic controls and a verified native-competence audit.
\begin{abstract}
Does repetition erase information from recurrent state, or merely change how that information is represented and read out? Low covariance dimension, contracting tangent dynamics, and poor task performance do not resolve this question individually. We formulate an experiment in which a fixed-weight network writes a cue once, processes either a periodic distractor body or an equal-length matched control, and then answers a query. We organize analytic nulls and counterexamples that constrain its interpretation. For fixed affine dynamics, cue differences are independent of body order at equal horizon; commuting switched-affine dynamics obey an analogous null under matched symbol counts. Nevertheless, order-dependent offsets can change native error without erasing the cue. Conversely, locally contracting dynamics can preserve categorical memories in distinct basins, while isotropic attenuation can reduce noise-resolved memory without changing normalized dimension. These controls motivate separate measurements of conditioned geometry, complete-state propagation, and task information. An audit of four single-seed GRU/LSTM pilots illustrates the prerequisite: all fail short-horizon native competence, making repetition-induced forgetting inconclusive. The contribution is a falsification framework and a controlled experimental design for recurrent-memory claims; the current evidence does not establish a learned collapse mechanism.
\end{abstract}
